\scanpage{41}\mspage{3}

\deleted{\unsure{c}} \uline{Remarque} Le groupe des automorphismes de la situation donnée se dévisse en $\unsure{\mathfrak S|\mathscr S}_I$ ($\simeq\mathfrak S_3$) \deleted{$\operatorname{Aut}G_j\times_{\unsure{\operatorname{Aut}\operatorname{Ext}G_i}}\operatorname{Aut}G_i$} et

\[
\underset{\simeq\Gamma\ \text{si les }G_i\ \text{\unsure{comm.}}}{\underbrace{\prod_{i\in I}{}_{\Gamma}\operatorname{Aut}(G_{i})}}\,,
\quad\text{où }\Gamma^{\deleted{\unsure{\varepsilon}}}=
 \underbrace{\operatorname{Aut}\operatorname{Ext}(G_i)}_{\substack{\text{indépendant de }i\\\text{mod syst. tran.}\\\text{d'iso}}}
\]

$\prod G_i$, opérant sur $(V_i)$ via

\[
g_jx_i g_k^{-1},\quad g_kx_jg_i^{-1},\quad g_ix_kg_j^{-1}
\]

% slanted
\margin{\unsure{oubli !|voir !|restriction|produit} au vrai \unsure{compte|conjugués}, il faut prendre \deleted{$\unsure{G^{I}/Z}\ \unsure{Z^I/\mathrm{diag}|Z^I/\mathrm{diagonale}}$} où $Z$ est le ``centre commun'' des $G_i$}

\margin{Il est plus joli de dévisser en $\mathfrak S_I$, \deleted{Aut}$\operatorname{Ext}(G_{\unsure{i|*}})$, et $G^I/Z$ (envoyé diagonalmt en $Z^I\subset G^I$}

\textcircled{c} Soit maint$^{\mathrm{nt}}$ $I$ \unsure{un|ens} à \uline{quatre} éléments, et \deleted{\unsure{d'un}} \unsure{se|on|on se} \unsure{proposons|propose} \unsure{nous|encore} d'étudier les

\[
 V\xrightarrow{\varphi_i}V_i\qquad i\in I
\]

dans $C$, tels que $\forall J\in I$, $\operatorname{card} J=3$, on ait

\[
V\xrightarrow{\varphi_J}V_J\qquad\text{iso.}
\]

Pour toute \deleted{subdivision} numérotation

\[
 I=\{i_1,i_2,i_3,i_4\}
\]

on va définir une flèche

\[
 u_{(i_1,i_2),(i_3,i_4)}:V_{i_1}\xrightarrow{\sim}V_{i_2}
 \tag*{(1)}
\]

\margin{\unsure{24|Il} \unsure{quatrts|quantités|quadruplets|par}}

C'est l'iso can. défini en a), quand on fait abstraction de $\varphi_{i_4}$, \deleted{\unsure{Il|donc}}

% archive sheet number: 40
